\scanpage{3}
\margin{Juillet--Août \unsure{1967}}

Anneaux \deleted{\unsure{semi-}}simpliciaux canoniques et modules
stratifiés. \deleted{\unsure{cristaux}}
[Formellement, les développements qui suivent sont indépendants
des \deleted{\unsure{crist.}} histoires de Sites cristallins ou
\unsure{Antipiants|stratifiants}].

1. Soit $X\longrightarrow S$ un préschéma relatif.
Pour tout entier $n\geq0$, soit
\[
X^{[n]}=\underbrace{X\times_S\cdots\times_S X}_{n+1}
\]
considéré comme objet \deleted{de $\unsure{X/(\mathrm{Sch}/S)}$}
de $(\mathrm{Sch})_{/S}$, muni de \deleted{\unsure{diag}}
l'immersion
$X\xrightarrow{\operatorname{diag}_n}X^{[n]}$.
Soit $X^{\widehat{[n]}}$ son complété formel le long de $X$.
Pour $n\geq0$ variable, on trouve un objet simplicial à valeurs
dans $\mathrm{Schf}/X$, les espaces sous-jacents aux
$X^{\widehat{[n]}}$ étant d'ailleurs tous égaux à $\operatorname{esp}(X)$.

Soit $\mathcal P^n$
$\deleted{\unsure{\mathcal O_{X^{\widehat{[n]}}}}}$
le faisceau structural de $X^{\widehat{[n]}}$,
\deleted{\unsure{vu dans}} considéré comme
\deleted{\unsure{anneau}}
\inserted{$\mathbb N$-pro-objet de \unsure{la}}
la catégorie des faisceaux d'anneaux sur $\operatorname{esp}(X)$,
i.e. dans $\operatorname{Pro}(\operatorname{Ann}(\operatorname{esp}(X)))$.
Pour $n$ variable, on trouve ainsi un Anneau cosimplicial
$\mathcal P^*$ sur $\operatorname{esp}(X)$, avec
$\mathcal P^0=\underline{\mathcal O}_X$.
Faire attention que ce \uline{n'est}
\margin{2}
